Foundation models for 3D computed tomography (CT) are pretrained to be adapted to clinical datasets, which usually carry several labels per scan, such as multiple findings in a classification cohort or multiple structures in a segmentation dataset. However, their published evaluations mostly keep the encoder frozen and train a lightweight head on its features, and therefore do not inform two decisions that adapting such a model requires: which parameters to train, and whether to train one model per label or a single model for the whole dataset. To address this, we compare \numClsMethodsN{} classification and \numSegMethodsN{} segmentation adaptation methods, ranging from frozen probes to low-rank adapters and full fine-tuning, each trained both per label and jointly with one default setting and a fixed training budget. We evaluate three public 3D CT encoders, one of which is restricted to head CT, on \numLabelsN{} binary labels from five public classification cohorts and \numStructuresN{} structures from five public segmentation datasets, totaling \numRunsMainN{} training runs. In classification, adapting the encoder with LoRA improves held-out AUROC by \magn{\numCiClsloralinAllMeandiff} over the linear probe and by \magn{\numCiClsloraattnAllMeandiff} over the frozen attentive probe, and fine-tuning all attention weights performs comparably at similar memory and time per training step, although LoRA updates \numParamsratioPillarZAttnftLora{} to \numParamsratioSpectreAttnftLora{} times fewer encoder weights. A single joint model per cohort is not detectably different from per-label models on average (AUROC difference $\numCiClsjointAllMeandiff$, 95\% interval $\numCiClsjointAllLodiff$ to $\numCiClsjointAllHidiff$). In segmentation, a multi-depth decoder on the frozen encoder improves Dice over a linear probe by \magn{\numSegdeltaTapctLinDecMeandiff} (TAP-CT) and \magn{\numSegdeltaSpectreLinDecMeandiff} (SPECTRE), most of the gain of adapting the encoder, whereas per-structure models remain more accurate than joint models at our training budgets. Based on these findings, we provide practical ``rule-of-thumb'' guidelines for adapting 3D CT foundation models. Our code and data splits are available at \url{https://github.com/georgeyiasemis/beyond_frozen_encoders}.
